% Generated draft: requires mathematical and source review.
\documentclass[11pt]{article}
\usepackage[margin=25mm]{geometry}
\usepackage{fontspec}
\setmainfont{Latin Modern Roman}
\usepackage{amsmath,amssymb,mathtools,unicode-math}
\setmathfont{Latin Modern Math}
\usepackage{enumitem,needspace,xurl,hyperref}
\setlength{\parindent}{0pt}
\setlength{\parskip}{5pt}
\setlength{\emergencystretch}{3em}
\newcommand{\N}{\mathbb{N}}
\newcommand{\Z}{\mathbb{Z}}
\newcommand{\Q}{\mathbb{Q}}
\newcommand{\R}{\mathbb{R}}
\newcommand{\C}{\mathbb{C}}
\newcommand{\F}{\mathbb{F}}
\newcommand{\A}{\mathbb{A}}
\newcommand{\PP}{\mathbb{P}}
\newcommand{\E}{\mathbb{E}}
\newcommand{\eps}{\varepsilon}
\newcommand{\dd}{\,\mathrm{d}}
\begin{document}
\section*{UnsolvedMath ID: 10400041}
\textbf{Source problem code:} catalogue entry
\subsection{Definitions and mathematical statement}
For an integer \(\mu>2\), a \(\mu\)-component link is an embedding of \(\mu\) disjoint circles in \(S^3\), considered up to ambient isotopy. Component ordering and orientation must be fixed when stating a classification; the excerpt does not specify them.

A \(\Delta\)-move is the standard local three-strand replacement associated with a Borromean-type tangle, supported in a ball and preserving its boundary endpoints. The source uses the technical term ``\(\Delta\) link homotopic'' without reproducing the allowed component pattern for the three strands or the equivalence-generating moves. Unrestricted \(\Delta\)-equivalence, moves entirely on one component, other component restrictions, and ordinary link homotopy by self-crossing changes must not be identified silently.

With the source's equivalence relation \(\sim_\Delta\) fixed, find necessary and sufficient conditions for
\[
 L\sim_\Delta L'
\]
for arbitrary \(\mu\)-component links with \(\mu\ge3\). Equivalently, provide a complete criterion or invariant \(I\) satisfying \(L\sim_\Delta L'\) if and only if \(I(L)=I(L')\), with the intended label/orientation conventions. The source does not require a finite list or a particular algorithm, and neither is imposed.

The supplied summary says the 2002 question was not obviously resolved but gives no classification theorem or modern citation. That open/triage wording is a qualified source assessment, not independently verified current status. The missing move convention is essential before any candidate criterion can be evaluated.
\subsection{Short English statement}
Classify links with more than two components up to the source's \(\Delta\)-link-homotopy relation. The precise component restrictions on the moves, and orientation/label conventions, are missing and must be recovered before stating a complete criterion.
\subsection{Sources}
\begin{itemize}
\item Local registry: \url{data/unsolvedmath/problems.json}, permanent UnsolvedMath ID \(10400041\), source code catalogue entry; supplied statement, background, and research summary, updated 2026-08-06.
\item Ohtsuki, \emph{Problems on invariants of knots and 3-manifolds} (2002), Problem 2.19, attributed to Y. Ohyama, PDF page 43: \url{https://msp.org/gtm/2002/04/gtm-2002-04-024s.pdf}. Supplied PDF-font-structure extraction accessed 2026-07-29; provenance distinguishes 231 numbered items.
\item No exact modern classification or separate move-definition reference is supplied.
\end{itemize}
\end{document}
